\chapter{Teorie ohledně obsahu aplikace}
\label{chap:content}

V této kapitole jsou popsány teoretické základy týkající se obsahu aplikace. Nejprve je uveden přehled různých datových struktur, které aplikace podporuje, včetně jejich vlastností, operací a požadavků na vizualizaci. Dále jsou diskutovány třídící algoritmy, které budou implementovány, a jejich význam pro pochopení základních konceptů informatiky. Nakonec je popsán přístup k vizualizaci těchto datových struktur a algoritmů, včetně použitých technik a metod pro efektivní zobrazení a interakci s uživatelem.~\cite{levitin}

\todo{Popsat teorii k obsahu aplikace, vycházet z Levitina a jiných zdrojů jestli nebude stačit. (Ideálně asi obrázek ke každé struktuře, ale zas to nerozepisovat moc do detailu, spíš jen základní vlastnosti a operace. U algoritmů pak popsat principy fungování a složitosti, ale tam toho fakt nebude moc.)}

\section{Datové struktury}
\label{sec:data-structures}

V této sekci jsou popsány datové struktury, které budou vizualizovány v aplikaci. Každá datová struktura je charakterizována svými vlastnostmi, operacemi a požadavky na vizualizaci.

\section{Třídící algoritmy}
\label{sec:algorithms}

V této sekci jsou popsány třídící algoritmy, které budou implementovány v aplikaci. Každý algoritmus je charakterizován svým principem fungování, složitostí a významem pro pochopení základních konceptů informatiky.
